\documentclass[11pt,a4paper]{article}
\usepackage[margin=25mm,headheight=14pt]{geometry}
\usepackage{fontspec}
\setmainfont{lmroman10-regular.otf}[
  BoldFont=lmroman10-bold.otf,
  ItalicFont=lmroman10-italic.otf,
  BoldItalicFont=lmroman10-bolditalic.otf]
\setsansfont{lmsans10-regular.otf}[BoldFont=lmsans10-bold.otf]
\setmonofont{lmmono10-regular.otf}
\usepackage{amsmath,amssymb,amsthm,mathtools}
\usepackage{booktabs,array,enumitem,microtype,xurl}
\usepackage[colorlinks=true,linkcolor=black,citecolor=black,urlcolor=blue]{hyperref}
\usepackage{fancyhdr}
\pagestyle{fancy}
\fancyhf{}
\fancyhead[L]{\small Directional differentiation and interpolation}
\fancyhead[R]{\small Colombo's determinant problem}
\fancyfoot[C]{\thepage}
\renewcommand{\headrulewidth}{0.3pt}
\setlength{\parindent}{1.2em}
\setlength{\parskip}{0pt}
\setlength{\emergencystretch}{2em}
\setlist{topsep=4pt,itemsep=3pt}
\numberwithin{equation}{section}
\newtheorem{theorem}{Theorem}[section]
\newtheorem{lemma}[theorem]{Lemma}
\newtheorem{proposition}[theorem]{Proposition}
\theoremstyle{remark}
\newtheorem{remark}[theorem]{Remark}
\newcommand{\R}{\mathbb R}
\newcommand{\N}{\mathbb N}
\newcommand{\RP}{\mathbb {RP}^{1}}
\newcommand{\supp}{\operatorname{supp}}
\newcommand{\rank}{\operatorname{rank}}
\newcommand{\Pf}{\operatorname{Pf}}
\newcommand{\pdeg}{\operatorname{deg}}
\newcommand{\lean}[1]{\texttt{\detokenize{#1}}}
\newcommand{\dd}{\mathcal D}
\hypersetup{
  pdftitle={Directional differentiation and interpolation for Colombo's determinant problem},
  pdfauthor={HappyAny},
  pdfsubject={A complete determinant criterion with a single-file Lean proof},
  pdfkeywords={difference-power matrix, projective roots, Rolle theorem, interpolation, Lean}
}
\title{\bfseries Directional differentiation and interpolation\\
for Colombo's determinant problem}
\author{HappyAny}
\date{6 September 2026 \quad\textnormal{\small Working draft}}

\begin{document}
\maketitle
\thispagestyle{plain}

\begin{abstract}
For distinct real numbers $x_1,\ldots,x_N$, we give a self-contained proof that
the matrix $[(x_j-x_i)^D]$ is nonsingular exactly when $N\le D+1$ and at least
one of $N,D$ is even. The proof is organized around an auxiliary polynomial
whose values have the signs of a putative kernel vector. Its real roots select
directional derivatives that transform the associated homogeneous power sum
into a strictly positive sum of even powers. A projective Rolle inequality
then bounds the original root count. Ordinary interpolation treats opposite
parities. When both parameters are even, a skew-symmetric derivative pairing
annihilates the highest interpolation coefficient and lowers the bound by two.
This construction gives a common proof framework for all nonsingular cases
and closely follows an accompanying single-file Lean formalization. We explain
its relation to Ma's Pfaffian positivity method and to the algebraic root-count
method of Li, Tie, Wang, and Liu, including the common parity obstruction in
the odd-exponent case.
\end{abstract}

\noindent\textbf{Keywords.} Difference-power matrix; real projective roots;
directional derivative; Lagrange interpolation; formal verification.

\section{The criterion and the proof architecture}

For integers $N\ge2$, $D\ge1$, and pairwise distinct real nodes
$x=(x_1,\ldots,x_N)$, put
\begin{equation}\label{eq:matrix}
 A_{N,D}(x)=\bigl[(x_j-x_i)^D\bigr]_{i,j=1}^{N}.
\end{equation}
We prove the following complete criterion.

\begin{theorem}\label{thm:main}
For the matrix in \eqref{eq:matrix},
\begin{equation}\label{eq:criterion}
 \det A_{N,D}(x)\ne0
 \quad\Longleftrightarrow\quad
 N\le D+1\ \text{and}\ (N\text{ is even or }D\text{ is even}).
\end{equation}
\end{theorem}

The criterion is already stated in Ma's treatment of Colombo's problem
\cite{Ma}. Li, Tie, Wang, and Liu give an algebraic proof of the remaining
odd-exponent branch for even matrix size \cite{Li}. Both papers combine that
branch with known results for even exponents. The present paper concerns a
different organization of the proof, particularly the construction that
closes the even-size, even-exponent case. The numerical criterion itself is
not presented as a new result.

Suppose that $A_{N,D}(x)c=0$ for a nonzero vector $c$. The associated polynomial
\begin{equation}\label{eq:kernel-polynomial}
 f(t)=\sum_{j=1}^{N}c_j(t-x_j)^D
\end{equation}
vanishes at every node. We will bound its real roots by constructing a
polynomial $g$ with
\begin{equation}\label{eq:alignment}
 c_jg(x_j)>0\qquad\text{whenever }c_j\ne0.
\end{equation}
The basic estimate is that $\deg g\le m\le D$ and $D-m$ even imply at most
$m$ projective real roots, provided at least two distinct nodes occur in the
support. The same estimate establishes that $f$ is nonzero, so no separate
Vandermonde independence argument is needed in the sufficiency proof.

There are two applications. Interpolating $g(x_j)=c_j$ gives $\deg g\le N-1$
and settles opposite parities. In the remaining case, both $N$ and $D$ are
even. The initial root bound forces a factorization $f=Pq$, where
$P(t)=\prod_j(t-x_j)$ and $q$ has no real root. A derivative pairing then
produces an aligned interpolation polynomial of degree at most $N-2$.
This last reduction is the main construction emphasized here.

\section{Projective root counts and directional Rolle}

For a nonzero real polynomial $p$ with $\deg p\le d$, let
\begin{equation}\label{eq:root-count}
 R_d(p)=Z_\R(p)+d-\deg p,
\end{equation}
where $Z_\R(p)$ counts real roots with multiplicity. Equivalently, $R_d(p)$
counts the zeros of the degree-$d$ homogenization
\[
 p^{[d]}(u,v)=v^dp(u/v)
\]
on $\RP$, with multiplicity. The displayed expression is understood as a
homogeneous polynomial, including at $v=0$. The term $d-\deg p$ is exactly
the multiplicity at $[1:0]$. A nonzero form of degree zero has root count zero.
For a nonzero homogeneous binary form $F$, we write $R(F)$ for this count.

For $w=(a,b)\ne(0,0)$, define
\[
 \dd_wF=a\,\partial_uF+b\,\partial_vF.
\]

\begin{lemma}[Projective Rolle inequality]\label{lem:rolle}
Let $F$ be a real homogeneous binary form of degree $d\ge1$. If
$\dd_wF\ne0$, then $F\ne0$ and
\begin{equation}\label{eq:rolle}
 R(F)\le R(\dd_wF)+1.
\end{equation}
\end{lemma}

\begin{proof}
Nonvanishing of the derivative implies $F\ne0$. An invertible real linear
change of coordinates preserves projective root multiplicities and can send
the direction of differentiation to $\partial_u$. In those coordinates,
write
\[
 F(u,v)=v^\ell G(u,v),\qquad G(1,0)\ne0,
 \qquad g(t)=G(t,1).
\]
Then $\deg g=d-\ell$. The derivative is nonzero only if $d-\ell\ge1$.
Its factor at infinity has the same multiplicity $\ell$, since
$\partial_uG(1,0)=(d-\ell)G(1,0)\ne0$. Thus
\[
 R(F)=\ell+Z_\R(g),\qquad
 R(\partial_uF)=\ell+Z_\R(g').
\]
The ordinary Rolle theorem, with multiplicities, gives
$Z_\R(g')\ge Z_\R(g)-1$. Indeed, a root of multiplicity $r$ contributes
at least $r-1$ to the derivative, and between any two successive distinct
roots there is an additional derivative root. If $g$ has no real root, the
inequality is immediate. This proves \eqref{eq:rolle}.
\end{proof}

\begin{lemma}[A positive terminal derivative]\label{lem:terminal}
Suppose that $k$ successive nonzero real directional derivatives transform
a degree-$d$ form $F$ into a nonzero scalar multiple of a form that is
strictly positive on $\R^2\setminus\{0\}$. Then $F\ne0$ and $R(F)\le k$.
\end{lemma}

\begin{proof}
The terminal form has no projective real zero. Its nonvanishing implies
nonvanishing of every preceding form. Apply Lemma~\ref{lem:rolle} backwards
through the $k$ derivatives. The case $k=0$ is included.
\end{proof}

The directional derivative needed below has the simple form
\begin{equation}\label{eq:directional-power}
 \dd_{(a,1)}(u-xv)^d=d(a-x)(u-xv)^{d-1}.
\end{equation}
It lowers the exponent while multiplying the coefficient by a linear
function of the node. Choosing that linear function is the essential degree
of freedom in the argument.

\section{A root bound from an aligned polynomial}

\begin{theorem}[Aligned polynomial root bound]\label{thm:aligned}
Let $x_1,\ldots,x_N$ be distinct real numbers, and let $c\in\R^N$ have at
least two nonzero coordinates. Let $0\le m\le d$ with $d-m$ even. Suppose
that a real polynomial $g$ satisfies $\deg g\le m$ and
\[
 c_jg(x_j)>0\qquad(j\in\supp c).
\]
Then the polynomial
$f(t)=\sum_jc_j(t-x_j)^d$ is nonzero, and
\begin{equation}\label{eq:aligned-bound}
 R_d(f)\le m.
\end{equation}
\end{theorem}

\begin{proof}
The sign condition implies $g\ne0$. Remove all real roots of $g$, counting
multiplicities, to write
\[
 g(t)=Q(t)\prod_{\nu=1}^{\ell}(t-a_\nu),
 \qquad Q(t)\ne0\quad(t\in\R),
 \qquad \ell\le\deg g\le m.
\]
By continuity, $Q$ has constant sign. Choose $\sigma\in\{1,-1\}$ with that
sign and put $h_0(t)=\sigma\prod_\nu(t-a_\nu)$. Then
\[
 c_jh_0(x_j)=c_jg(x_j)\,\frac{\sigma}{Q(x_j)}>0
 \qquad(j\in\supp c).
\]
If $d-\ell$ is even, set $h=h_0$ and $k=\ell$. Otherwise choose
$b>\max_jx_j$ and set
\[
 h(t)=(b-t)h_0(t),\qquad k=\ell+1.
\]
In the latter case $\ell<m$, because $\ell=m$ would contradict the
assumption that $d-m$ is even. In either case we have
\begin{equation}\label{eq:split-aligned}
 h(t)=\kappa\prod_{\nu=1}^{k}(t-b_\nu),\quad
 k\le m,\quad d-k\text{ even},\quad
 c_jh(x_j)>0\ (j\in\supp c),
\end{equation}
where $\kappa\ne0$ and all $b_\nu$ are real. Repetitions among the $b_\nu$
are allowed. No $b_\nu$ is an active node, by the strict sign condition.

Consider the homogeneous form
\[
 F(u,v)=\sum_jc_j(u-x_jv)^d.
\]
Repeated use of \eqref{eq:directional-power} gives
\begin{align}\label{eq:terminal-form}
 \dd_{(b_k,1)}\cdots\dd_{(b_1,1)}F
 &=\frac{d!}{(d-k)!}\sum_jc_j
        \prod_{\nu=1}^{k}(b_\nu-x_j)(u-x_jv)^{d-k}\notag\\
 &=\frac{d!}{(d-k)!}\frac{(-1)^k}{\kappa}
        \sum_jc_jh(x_j)(u-x_jv)^{d-k}.
\end{align}
The last sum is strictly positive at every nonzero $(u,v)$. Its nonzero
coefficients are positive and its exponent is even. If the exponent is
positive, two linear forms $u-x_iv$ and $u-x_jv$ with distinct nodes cannot
both vanish at a nonzero point. If the exponent is zero, the sum is a
positive constant. Lemma~\ref{lem:terminal} now gives $F\ne0$ and
$R(F)\le k\le m$. Dehomogenization identifies this statement with
\eqref{eq:aligned-bound}.
\end{proof}

\begin{remark}\label{rem:role}
The theorem uses the degree of a polynomial encoding the coefficient signs,
not merely the number of summands. Its terminal positivity proves
nonvanishing as well as a root bound. It is a formulation of a classical
sign-control mechanism, closely related to the root-count perspective of
Sylvester and Reznick \cite{Reznick}; no claim of a new general rule of signs
is needed for its use here.
\end{remark}

\section{Proof of the complete determinant criterion}

\subsection{The two necessary conditions}

Let $V$ be the $N\times(D+1)$ matrix with $V_{ik}=x_i^k$, indexed by
$0\le k\le D$. Let $C$ be the $(D+1)\times(D+1)$ matrix whose only nonzero
entries are
\[
 C_{k,D-k}=(-1)^k\binom Dk.
\]
The binomial theorem gives
\begin{equation}\label{eq:factorization}
 A_{N,D}(x)=VCV^{\mathsf T},\qquad \rank A_{N,D}(x)\le D+1.
\end{equation}
Thus $N>D+1$ forces singularity. Also
\[
 A_{N,D}(x)^{\mathsf T}=(-1)^D A_{N,D}(x).
\]
If $N$ and $D$ are both odd, taking determinants gives
$\det A=-\det A$, hence $\det A=0$. These are precisely the two obstructions
on the right side of \eqref{eq:criterion}.

\subsection{The kernel polynomial and its support}

For the rest of the proof assume $N\le D+1$, and suppose, for a
contradiction, that $Ac=0$ for some $c\ne0$. The support of $c$ contains
at least two indices. Otherwise, if only $c_j$ is nonzero, any row $i\ne j$
would give
\[
 (Ac)_i=c_j(x_j-x_i)^D\ne0,
\]
contradicting the kernel equation.
For the polynomial \eqref{eq:kernel-polynomial},
\begin{equation}\label{eq:node-roots}
 f(x_i)=(-1)^D(Ac)_i=0\qquad(1\le i\le N).
\end{equation}
Consequently, once $f\ne0$ has been established, $R_D(f)\ge N$.

\subsection{Opposite parities}

Suppose that $N$ and $D$ have opposite parity. Let $g$ be the Lagrange
interpolation polynomial satisfying
\[
 \deg g\le N-1,\qquad g(x_i)=c_i.
\]
At every active node, $c_ig(x_i)=c_i^2>0$. Apply
Theorem~\ref{thm:aligned} with $m=N-1$. Its hypotheses hold because
$N-1\le D$ and $D-(N-1)$ is even. It follows that $f\ne0$ and
\[
 N\le R_D(f)\le N-1,
\]
a contradiction. This argument includes the threshold $D=N-1$ and both
opposite-parity choices.

\subsection{Even size and even exponent: interpolation loses a degree}

Suppose now that $N$ and $D$ are even. Since $N\le D+1$, parity gives
$N\le D$. Start with the same interpolation polynomial $g_0(x_i)=c_i$,
but apply Theorem~\ref{thm:aligned} with $m=N$. We obtain $f\ne0$ and
\[
 N\le Z_\R(f)\le R_D(f)\le N.
\]
All inequalities are equalities. In particular, $\deg f=D$, the nodes are
all the real roots of $f$, and each is simple. We can therefore write
\begin{equation}\label{eq:rootless-factor}
 f=Pq,\qquad P(t)=\prod_{i=1}^{N}(t-x_i),
 \qquad q(t)\ne0\quad(t\in\R).
\end{equation}
Set
\begin{equation}\label{eq:weights}
 w_i=P'(x_i)=\prod_{j\ne i}(x_i-x_j)\ne0.
\end{equation}
Differentiating \eqref{eq:rootless-factor} at a node gives
\begin{equation}\label{eq:derivative-values}
 f'(x_i)=w_iq(x_i).
\end{equation}

The exponent $D-1$ is odd. Exchanging the indices in a finite double sum
therefore proves the identity
\begin{equation}\label{eq:skew-pairing}
 \sum_{i=1}^{N}c_if'(x_i)
 =D\sum_{i,j=1}^{N}c_ic_j(x_i-x_j)^{D-1}=0.
\end{equation}
This identity is valid for every vector $c$; the kernel equation is not
needed for the cancellation itself.

Interpolate the new data
\begin{equation}\label{eq:new-data}
 g(x_i)=c_if'(x_i)w_i,\qquad \deg g\le N-1.
\end{equation}
The Lagrange formula reads
\[
 g(t)=\sum_i g(x_i)\frac{P(t)}{(t-x_i)w_i},
\]
where each quotient $P(t)/(t-x_i)$ is a monic polynomial of degree $N-1$.
Taking the coefficient of $t^{N-1}$ and using \eqref{eq:skew-pairing},
\begin{equation}\label{eq:degree-drop}
 [t^{N-1}]g=\sum_i\frac{g(x_i)}{w_i}
 =\sum_i c_if'(x_i)=0.
\end{equation}
Hence $\deg g\le N-2$, with the usual convention for the zero polynomial.

Since $q$ has no real root, $q(0)q(x_i)>0$ for every $i$. The scaled
polynomial $h=q(0)g$ has degree at most $N-2$, and
\begin{align}\label{eq:new-alignment}
 c_ih(x_i)
 &=q(0)c_i^2f'(x_i)w_i\notag\\
 &=q(0)q(x_i)(c_iw_i)^2>0\qquad(i\in\supp c).
\end{align}
Theorem~\ref{thm:aligned} applies again, now with $m=N-2$, because
$D-(N-2)$ is even. It gives
\[
 N\le R_D(f)\le N-2,
\]
the final contradiction. No ordering of the nodes is used in this
construction. The case $N=2$ is included: the last auxiliary polynomial
then has degree at most zero.

\begin{proof}[Completion of Theorem~\ref{thm:main}]
Equation~\eqref{eq:factorization} and odd-dimensional skew-symmetry give
necessity. Under the stated conditions, the opposite-parity and even-even
arguments exclude every nonzero kernel vector. A square matrix over $\R$
with trivial kernel has nonzero determinant, proving sufficiency.
\end{proof}

\section{Comparison with the two existing approaches}

\subsection{Ma: positivity of a Pfaffian integral}

Ma's odd-exponent proof \cite{Ma} constructs a strict sign for the Pfaffian.
For ordered $N=2r$ nodes, its output is
\[
 (-1)^{\binom r2}\Pf(A_{N,D}(x))>0.
\]
The proof passes through paired determinants of truncated powers, B-spline
and Marsden factorizations, total nonnegativity, Volterra propagation, and a
Beta--de Bruijn integral representation. Positivity on an explicit region
then gives a strictly positive integral.

The proof above excludes a kernel vector through a polynomial root bound.
Its positive object is the terminal even-power sum in
\eqref{eq:terminal-form}. This avoids constructing the paired determinants
and integral representation. Ma's formalization also establishes intermediate
positivity statements and the explicit Pfaffian sign \cite{MaCode}; those
are additional outputs beyond the nonvanishing interface formalized here.

\subsection{Li, Tie, Wang, and Liu: a closely related root-count contradiction}

Li et al.\ \cite{Li} associate a real binary form to a kernel vector and
apply the Sylvester--Reznick bound relating real projective factors to the
length of a power-sum representation. For $N$ even and $D\ge N+1$ odd,
their argument has the shape
\[
 N+1\le R_D(f)\le |\supp c|\le N.
\]
Our corresponding argument has the shape $N\le R_D(f)\le N-1$. These
inequalities express the same parity obstruction: for nonzero $f$ of degree
at most $D$,
\[
 R_D(f)\equiv D\pmod2,
\]
because nonreal roots occur in conjugate pairs. When $D$ is odd and $N$ is
even, either side of the root-count comparison can incorporate this parity
improvement. The change from $N$ to $N-1$ alone does not distinguish a new
mathematical obstruction.

The relevant difference lies in the construction of the bound. Their Lean
proof of the required finite-sum version \cite{LiCode} translates one node
to zero, reflects the polynomial, and differentiates. One term becomes
constant and disappears; induction reduces the number of summands to two.
In Theorem~\ref{thm:aligned}, the chosen directions instead encode signs,
and differentiation ends at a positive sum of even powers. Both developments
use one-variable real polynomials with a degree parameter to account for
roots at infinity.

\subsection{The even-even construction and formal scope}

The papers \cite{Ma,Li} cite existing even-exponent results. Their source
repositories have different scopes. Ma's checked repository formalizes the
odd branch. The checked version of Li et al.'s repository also contains
the even branch \cite{LiCode}. In its even-even proof, the first root bound
likewise yields \eqref{eq:rootless-factor}. It then establishes strict sign
alternation of kernel vectors of a reduced odd-order skew matrix, using an
explicit tangent-node configuration and a continuous deformation. A polar
derivative identity yields a contradictory dot product.

Equations~\eqref{eq:skew-pairing}--\eqref{eq:new-alignment} provide the
alternative used here: the skew identity cancels a leading interpolation
coefficient, which directly improves the root bound. This is a concrete
simplification of the auxiliary construction required for this criterion.
It does not include the separate alternating-kernel theorem proved in
\cite{LiCode}.

The relationship with prior work is consequently twofold. The odd branch
belongs to the same root-count family as \cite{Li,Reznick}, while the
aligned-polynomial formulation and the even-even interpolation construction
give a compact route through all parameter cases. No priority claim for
the determinant criterion or the underlying classical root-count principle
is made.

\section{The accompanying Lean proof}

The complete criterion is formalized in the single file
\lean{DifferencePowerSingle.lean}, using Lean 4.32.0 and mathlib v4.32.0.
It contains 958 lines, including comments and blank lines, and imports only
\lean{Mathlib}. All definitions and lemmas specific to this problem occur
in that file. Its final declaration is
\begin{quote}\small\ttfamily\raggedright
theorem determinant\_ne\_zero\_iff \{N D : $\N$\}\par
\quad (hN : $2\le N$) (\_hD : $1\le D$)\par
\quad (x : Fin N $\to\R$) (hx : Function.Injective x) :\par
\quad (Matrix.det (fun i j : Fin N => (x j - x i)\textasciicircum D)) $\ne0$\par
\quad $\leftrightarrow$ N $\le$ D+1 $\land$ (Even D $\lor$ Even N)
\end{quote}
This is the full biconditional for arbitrary parameters and injective real
node functions, rather than a statement restricted to examples.

\begin{center}
\small
\begin{tabular}{@{}p{0.37\linewidth}p{0.58\linewidth}@{}}
\toprule
Mathematical step & Principal Lean declaration \\
\midrule
Rank and parity obstructions & \lean{necessary} \\
Kernel equations give roots & \lean{kernel_gives_roots} \\
Projective Rolle estimate & \lean{projective_rolle} \\
Theorem~\ref{thm:aligned} & \lean{root_bound_of_sign_polynomial} \\
Opposite parities & \lean{kernel_zero_of_opposite_parity} \\
Equation~\eqref{eq:skew-pairing} & \lean{derivative_pairing_zero} \\
Even-even interpolation & \lean{kernel_zero_of_even_even} \\
Theorem~\ref{thm:main} & \lean{determinant_ne_zero_iff} \\
\bottomrule
\end{tabular}
\end{center}

The code represents a degree-$d$ form by a polynomial in $\R[X]$ and counts
its finite roots together with the degree deficit at infinity. Translation
and polynomial reflection implement the coordinate changes needed for
directional differentiation. The proof of the even-even case uses exactly
the interpolation data in \eqref{eq:new-data} and the leading-coefficient
cancellation in \eqref{eq:degree-drop}.

The file was copied to a fresh directory and compiled with a search path
containing only the pre-existing mathlib and dependency caches. Compilation
completed successfully. The transitive axiom audit of the final theorem
reported
\begin{center}
\lean{[propext, Classical.choice, Quot.sound]}.
\end{center}
There is no unfinished proof placeholder or additional mathematical axiom
in the development. The source, its verification log, and a SHA-256 manifest
are supplied with this manuscript. The source is unchanged from the
verification performed before the subsequent literature comparison.

The checked external source versions were
\lean{9f57b40c7c} for Ma \cite{MaCode} and \lean{b98c9e55c9} for Li et al.
\cite{LiCode}. Their repositories were inspected for the comparison above;
neither complete external project was rebuilt locally for this manuscript.
Our formal verification claim concerns the accompanying criterion proof,
not the bibliographical or workflow statements in this paper.

\paragraph{Code and reproducibility.}
The complete single-file Lean proof, pinned project configuration, original
verification record, and manuscript sources are available in the public
repository
\begin{center}
\url{https://github.com/HappyAny/colombo-determinant-lean}.
\end{center}
The repository also includes this PDF and instructions for checking the
proof with Lean and mathlib 4.32.0.

\begin{thebibliography}{9}
\small
\bibitem{Ma}
Q. Ma, \emph{Colombo's Determinant Problem}, arXiv:2609.00101v1, 2026.
\url{https://arxiv.org/abs/2609.00101}.
The arXiv record cites an earlier Zenodo deposit dated 18 August 2026,
\url{https://doi.org/10.5281/zenodo.21993537}.

\bibitem{Li}
K. Li, L. Tie, P. Wang, and Z. Liu,
\emph{An algebraic proof of Colombo's difference-power determinant conjecture},
arXiv:2608.28274v1, 2026.
\url{https://arxiv.org/abs/2608.28274}.

\bibitem{Reznick}
B. Reznick, \emph{On the length of binary forms},
arXiv:1007.5485, 2010, Section 3.
\url{https://arxiv.org/abs/1007.5485}.

\bibitem{MaCode}
Q. Ma, \emph{Colombo odd-exponent theorem in Lean 4},
source snapshot \lean{9f57b40c7c}, accessed 6 September 2026.
\url{https://github.com/hkjtsgmc79-boop/colombo-odd-lean}.

\bibitem{LiCode}
\emph{Colombo1928}, the Lean repository accompanying \cite{Li},
source snapshot \lean{b98c9e55c9}, accessed 6 September 2026.
\url{https://github.com/ttieli/Colombo1928}.
\end{thebibliography}

\section*{AI assistance and development time}

We developed the mathematical proof and its Lean formalization in an
interactive session with \textbf{GPT-6 Astra}, with reasoning effort set to
\textbf{max}. According to our session record, this original development
phase took \textbf{approximately 50 minutes}.
For this original solution-and-formalization phase, \textbf{Codex was the
sole harness}, \textbf{fast mode was disabled}, and \textbf{no special
orchestration skills were used}.
Our prompts were brief: solve the supplied problem without
online search, formalize the proof in Lean, and provide a complete
single-file derivation. We did not supply a detailed proof strategy in
those prompts; the ordinary Codex instructions and tools remained in use.
We carried out the original problem-solving stage without web searches,
using an existing local Lean/mathlib installation. We subsequently
consulted the papers and repositories cited above for comparison and
manuscript preparation. The reported
50-minute duration refers to the original proof-development phase and
excludes that subsequent literature work and the writing of this paper.

The scope of Astra's training data and prior knowledge is unknown to us,
including whether related mathematics or texts were present. Prohibiting
online search during the session does not establish the absence of such
prior exposure. The timing and workflow observations are therefore
\textbf{for reference only}, as an uncontrolled case report. This
qualification concerns provenance and efficiency claims; the mathematical
criterion is separately supported by the Lean kernel check.

The case raises an open question: as model capabilities improve, does the
\emph{marginal benefit of additional task-specific orchestration} decrease?
A stronger model may perform more planning and error correction within
the model, reducing the need for some external coordination. This single
session does not establish that trend. Codex still supplied tool execution,
context management, and access to verification feedback, while Lean checked
the resulting proof. A useful comparison would vary both model and harness
on the same tasks under matched resource budgets, and measure verified
completion and cost over repeated runs.

\textbf{In particular, the mathematical result and its complete Lean proof
were obtained under an explicit prohibition on online search during the
original solution and formalization.}

\end{document}
